\begin{corollary}[6.5.4]
\label{IV.2.s6.5.4-fr}
Sean $X$, $Y$ dos preesquemas localmente noetherianos y $f:X\to Y$
un morfismo plano.
\begin{enumerate}
\item[\textup{(i)}] Si $X$ es normal en un punto $x$, $Y$ es normal en el punto
$f(x)$.
\item[\textup{(ii)}] Si, para todo $y\in f(X)$, $f^{-1}(y)$ es un preesquema
normal y si $Y$ es normal en un punto $y\in f(X)$, entonces $X$ es un
preesquema normal en todo punto de $f^{-1}(y)$.
\end{enumerate}
\end{corollary}

Esto resulta inmediatamente del criterio de normalidad de Serre (5.8.6), de
(6.4.1) aplicado para $k=2$ y (6.5.3) aplicado para $k=1$.

\oldpage[IV-2]{144}

\begin{remarks}[6.5.5]
\EGAReaderTarget{ega4.c04.item.6-5-5}{IV.2.s6.5.5-fr}
\medskip\noindent
\begin{enumerate}
\item[(i)] Sean $A$, $B$ dos anillos locales noetherianos,
$\varphi:A\to B$ un homomorfismo local tal que $B$ sea un $A$-módulo
\emph{plano}. Sea $k$ el cuerpo residual de $A$, y supongamos que los dos anillos
$A$ y $B\otimes_Ak$ verifican la propiedad $(R_i)$ (5.8.2); entonces \emph{no se
sigue necesariamente} que $B$ verifique $(R_i)$, incluso en el caso
particular en que $i=0$ o $i=1$ y en que $B$ es el \emph{completado} $\widehat A$ de
$A$. En efecto, Nagata ha dado un ejemplo en el que $A$ es \emph{normal} (por tanto,
verifica $(R_1)$), pero $\widehat A$ ni siquiera es reducido (por tanto, no
verifica $(R_0)$) [30]. No se puede aplicar a este caso la proposición (6.5.3)
porque las fibras $f^{-1}(y)$ no verifican necesariamente la propiedad
$(R_i)$ en los puntos distintos del punto cerrado de
$Y=\operatorname{Spec}(A)$. Mostraremos, sin embargo, más adelante (7.8.3,
(v)) que tales fenómenos no se presentan para los anillos locales
noetherianos que se encuentran con mayor frecuencia en las aplicaciones.

\item[(ii)] La propiedad de que un preesquema sea \emph{íntegro} no se
comporta en absoluto como las propiedades que acabamos de examinar en este
número y en los precedentes: puede ocurrir que $f:X\to Y$ sea un morfismo
plano de tipo finito, que todas las fibras $f^{-1}(y)$ sean regulares (e
incluso geométricamente regulares (6.7.6)) y que $Y$ sea íntegro sin que $X$
sea siquiera localmente íntegro. Por ejemplo, sea $k$ un cuerpo
algebraicamente cerrado de característica $0$, y sea
$A$ el anillo íntegro $k[U,V]/(UV\setmin(U+V)^3)$ (cuyo espectro es, por tanto, una
«~cúbica con punto doble~»); $A$ no es integralmente cerrado, y si $u$, $v$
son las imágenes canónicas de $U$, $V$ en $A$, la clausura integral de $A$
es el anillo $C=A[t]$, con
$t=(u\setmin v)/(u+v)$, que verifica la ecuación $t^2=1+u+v$, de donde se obtiene
$u=\frac{1}{2}(t^3+t^2\setmin t)$,
$v=\frac{1}{2}(\setmin t^3+t^2+t)$
y, por consiguiente, $C=k[t]$, isomorfo al anillo de polinomios en una
indeterminada sobre $k$. Si $\mathfrak m=Au+Av$, ideal maximal de $A$
(correspondiente al «~punto doble~» de la cúbica), hay dos ideales maximales
$\mathfrak n_1$, $\mathfrak n_2$ de $C$, generados respectivamente por $t\setmin 1$
y $t+1$. Sea entonces $B$ el subanillo del producto $C\times C$ formado por
los pares de polinomios $(f,g)$ tales que $f(1)=g(\setmin 1)$ y $f(\setmin 1)=g(1)$
($\operatorname{Spec}(B)$ es el esquema obtenido «~pegando~» dos ejemplares
de $\operatorname{Spec}(C)$, estando «~pegado~» el punto $\mathfrak n_1$
(resp. $\mathfrak n_2$) de uno de los dos ejemplares al punto
$\mathfrak n_2$ (resp. $\mathfrak n_1$) del otro; cf. cap.~V, donde esta
operación será discutida de manera general). Hay, por tanto, dos ideales
maximales $\mathfrak r_1$, $\mathfrak r_2$ de $B$ por encima del ideal
maximal $\mathfrak m$ de $A$. Además, como el proceso de «~pegado~» conmuta
con la localización y el completado, se comprueba fácilmente que los
homomorfismos canónicos
$\widehat A_{\mathfrak m}\to\widehat B_{\mathfrak r_1}$ y
$\widehat A_{\mathfrak m}\to\widehat B_{\mathfrak r_2}$
son biyectivos y, por consiguiente (Bourbaki, \emph{Alg. comm.}, cap.~III,
\S~3, n\textsuperscript{o}~5, prop.~10), $B_{\mathfrak r_1}$ y
$B_{\mathfrak r_2}$ son $A_{\mathfrak m}$-módulos planos, y tienen además el
mismo cuerpo residual que $A_{\mathfrak m}$. Para cualquier otro ideal
maximal $\mathfrak p$ de $A$, es inmediato que hay dos ideales maximales
$\mathfrak q_1$, $\mathfrak q_2$ de $B$ por encima de $\mathfrak p$, y que
los homomorfismos $A_{\mathfrak p}\to B_{\mathfrak q_1}$ y
$A_{\mathfrak p}\to B_{\mathfrak q_2}$ son biyectivos. Se ve así que el
morfismo $\operatorname{Spec}(B)\to\operatorname{Spec}(A)$ es plano y
finito, y que todas sus fibras son geométricamente regulares (es incluso
\emph{\'etale}, como veremos más adelante (17.6.3)); sin embargo, es
inmediato que $B$ no es íntegro.
\end{enumerate}
\end{remarks}
